\documentclass[11pt]{article}
\usepackage[T1]{fontenc}\usepackage[utf8]{inputenc}
\usepackage{lmodern,amsmath,amssymb,amsthm,mathtools,microtype}
\usepackage[margin=1in]{geometry}
\usepackage[colorlinks=true,linkcolor=blue,urlcolor=blue]{hyperref}
\newtheorem{theorem}{Theorem}[section]\newtheorem{lemma}[theorem]{Lemma}
\newtheorem{corollary}[theorem]{Corollary}\newtheorem{proposition}[theorem]{Proposition}
\theoremstyle{remark}\newtheorem{remark}[theorem]{Remark}
\newcommand{\E}{\mathbb E}\newcommand{\PP}{\mathbb P}\newcommand{\R}{\mathbb R}
\setlength{\emergencystretch}{2em}
\hypersetup{pdftitle={Finite-width circuit learning below mean-field scaling},pdfauthor={Kaarel Hanni; Dmitry Vaintrob; OpenAI Codex}}
\title{Finite-width circuit learning\\below mean-field scaling}
\author{Kaarel H\"anni \and Dmitry Vaintrob \and OpenAI Codex\thanks{AI assistant; contributed to mathematical analysis and drafting.}}
\date{October 15, 2026}
\begin{document}\maketitle
\begin{abstract}
For fixed-depth Bayesian tanh networks with Gaussian scaling exponent
$\gamma<1$, we prove a finite-width sample bound of order
$\operatorname{poly}(d,C)N\log(eN)$, where $d$ is input dimension,
$C$ is the target's bounded-weight network complexity, and $N\ge C$
is the hidden width. The prior is restricted to a fixed positive output
margin on the entire Boolean cube, followed by exact noiseless conditioning
on the training labels. The guarantee concerns the actual posterior:
its correct-label probability exceeds $2/3$ on almost all test inputs.
Consequently, polynomial widths permit polynomial-sample learning of
parity, and every subexponential width sequence permits subexponential
sample complexity. These statements concern simultaneous growth of width
and dimension. They do not settle sample complexity when width first tends
to infinity at each fixed dimension. The proof gives an explicit Gaussian
probability bound by cloning a small target network and constraining only
linearly many parameter directions in width.
\end{abstract}

\section{Model, theorem, and order of limits}
Let $d\ge1$, $L\ge2$, and $\mathcal X=\{-1,1\}^d$.
A depth-$L$ network with common hidden width $N$ computes
\[
 h^0(x)=(1,x),\qquad
 h^s(x)=\tanh(W^sh^{s-1}(x))\quad(1\le s<L),
 \qquad z_w(x)=W^Lh^{L-1}(x).
\]
There are no additional biases, and the readout is linear. Fix a real
exponent $\gamma<1$. Under the product Gaussian law $\mu_{N,\gamma}$,
the independent centered weights have standard deviations
\begin{equation}\label{eq:scaling}
 \sigma_1=N^{(1-\gamma)/2}/\sqrt d,
 \qquad \sigma_s=N^{-\gamma/2}\quad(2\le s<L),
 \qquad \sigma_L=N^{-(\gamma+1)/2}.
\end{equation}
Fix an output margin $a>0$. For each Boolean function $g$, write
\[
 A_{N,g}^a=\{w:g(x)z_w(x)>a\ \text{for every }x\in\mathcal X\},
 \qquad
 \pi_N(g)=\frac{\mu_{N,\gamma}(A_{N,g}^a)}
 {\mu_{N,\gamma}(\bigcup_h A_{N,h}^a)}.
\]
Thus the global margin restriction precedes training. The posterior
$\Pi_N(\cdot\mid S_n)$ conditions this function prior on agreement with
noiseless examples $S_n=((x_i,f(x_i)))_{i=1}^n$, with independent $x_i\sim D$.
The distribution $D$ may be arbitrary on the cube.

Let $C=C_L(f)$ be the least integer $C\ge1$ for which a depth-$L$,
width-$C$ tanh network implements $f$ with strict margin $1/2$ and all
weights bounded in absolute value by $C$. We assume $C<\infty$.
The definition uses margin $1/2$ even when the Bayesian margin $a$ differs;
a constant readout rescaling handles this distinction.
Define
\begin{align*}
 R_N(S_n)&=\E_{g\sim\Pi_N}\PP_{x\sim D}\{g(x)\ne f(x)\},\\
 F_N(S_n)&=\PP_{x\sim D}\left\{
      \Pi_N(g(x)=f(x)\mid S_n)\le2/3\right\}.
\end{align*}
The equality case at $2/3$ is counted as a failure.

Here are explicit constants for the main bound:
\begin{align}
 b_L&=\sqrt{16\log9+8\log(4L)},&
 A&=2+3C^2+b_L,\notag\\
 D_0&=1+4a\sqrt2\,C^2(d+2)A^{L-2},&
 \rho_0&=\frac{a}{2D_0},\notag\\
 q&=d+2+(L-2)C,&
 B_0&=d(d+1)C^2+2(L-2)C^4+64a^2C^4.
 \label{eq:constants}
\end{align}
For every $N\ge C$, set
\begin{equation}\label{eq:cost}
 B_N=B_0N^\gamma
 +Nq\left[\log\frac{e\sqrt{2\pi}}{2\rho_0}
       +\frac{(L-1)(1-\gamma)}2\log N\right]+\log2.
\end{equation}

\begin{theorem}[Finite-width learning]\label{thm:learning}
For every $N\ge C$, the target has raw prior mass
\begin{equation}\label{eq:mass}
 \mu_{N,\gamma}(A_{N,f}^a)\ge e^{-B_N}.
\end{equation}
For every input distribution $D$ and every $\varepsilon,\delta\in(0,1)$,
if
\begin{equation}\label{eq:samples}
 n\ge\frac6\varepsilon
       \left[B_N+\log\frac6{\varepsilon\delta}\right],
\end{equation}
then, with probability at least $1-\delta$ over $S_n$,
\[
 R_N(S_n)\le\varepsilon/3,
 \qquad F_N(S_n)\le\varepsilon.
\]
At fixed $L,\gamma,a$, the cost obeys
\begin{equation}\label{eq:order}
 B_N\le p_{L,\gamma,a}(d,C)\,N\log(eN)
\end{equation}
for a polynomial $p_{L,\gamma,a}$ independent of width.
\end{theorem}

\begin{corollary}[Polynomial and subexponential widths]\label{cor:joint}
At fixed $L,\gamma,a,\varepsilon,\delta$, widths bounded by a polynomial
in $d,C_L(f)$ admit polynomial sample complexity under the actual posterior.
For full parity on $k$ bits, $C_L(f)=O_L(k\log(k+1))$.
Consequently, for every integer width sequence $N(k)\ge C_L(f)$ satisfying
$\log N(k)=o(k)$, there is a sample schedule $n(k)=\exp(o(k))$ for which
$F_{N(k)}(S_{n(k)})\le\varepsilon$ with probability at least $1-\delta$.
In particular, polynomial widths admit polynomial-sample parity learning.
\end{corollary}

The order of limits is part of the conclusion. The corollary permits width
to grow with dimension, for example $N(k)=\lceil e^{\sqrt{k}}\rceil$.
A different question fixes $k$ and its sample budget and then sends
$N\to\infty$. The sample bound~\eqref{eq:samples} grows with $N$, so it
does not establish an upper bound for that width-first limiting learner.
Nor does it contradict a lower bound proved for that learner.
An exponential-in-$k$ obstruction asserted uniformly over growing widths,
however, would conflict with Corollary~\ref{cor:joint}.

The guarantee is statistical. No efficient method for evaluating the
posterior is asserted. The following proofs are finite-dimensional and
do not require a Gaussian-process approximation.

\section{Cloning the target at an arbitrary width}
Choose a width-$C$ implementation of $f$ from its complexity definition.
Multiply its readout by $4a$. Write its hidden matrices as $B^s$ and
readout as $v$. The resulting teacher has strict margin greater than $2a$,
hidden entries bounded by $C$, and readout entries bounded by $4aC$.
Partition the $N$ neurons into $C$ groups with sizes $k_b$ differing by
at most one. Since $N\ge C$,
\begin{equation}\label{eq:groups}
 k_b\ge1,\qquad \frac{\max_b k_b}{\min_b k_b}\le2,
 \qquad \frac{N}{k_b}\le2C.
\end{equation}
The exact clone $w_*$ has first-layer row $B_b^1$ in group $b$, internal
entry $B^s_{ba}/k_a$ from group $a$ to group $b$, and readout entry
$v_b/k_b$ in group $b$. It computes precisely the teacher's output.
Each hidden activation is constant within every group.

Let $u_a=k_a^{-1/2}\mathbf1_{\text{group }a}$ and let $P$ be the
orthogonal projection onto their span. Restricting an internal cloned
matrix to this subspace gives entries $\sqrt{k_b/k_a}B^s_{ba}$.
Its Frobenius norm and~\eqref{eq:groups} imply
\begin{equation}\label{eq:clone-norms}
 \|W_*^s\|_{\mathrm{op}}\le\sqrt2\,C^2\quad(2\le s<L),
 \qquad \sqrt N\|W_*^L\|_2\le4a\sqrt2\,C^2.
\end{equation}
The cloned internal matrices vanish on the orthogonal complement of the
same group-constant subspace.

Put
\begin{equation}\label{eq:cylinder-radius}
 t=N^{(1-\gamma)/2}\ge1,
 \qquad \rho=\rho_0t^{-(L-1)}.
\end{equation}
The normalized coordinates $\sqrt d\,W^1_{ij}$,
$\sqrt N\,W^s_{ij}$ for $2\le s<L$, and $\alpha_j=NW^L_j$
are independent centered Gaussians with standard deviation $t$.
Impose the following cylinder event $\mathcal C_\rho$:

\begin{itemize}
\item Every $\sqrt d\,W^1_{ij}$ is within $\rho t$ of its cloned value.
\item For each internal row $j$ and group $a$, the coordinate
$\langle\sqrt N\,W^s_{j,\cdot},u_a\rangle$ is within $\rho t$ of
its cloned value $\sqrt{N/k_a}B^s_{ba}$, where $j$ belongs to group $b$.
\item Every $\alpha_j$ is within $\rho t$ of its cloned value $Nv_b/k_b$.
\end{itemize}

These are $Nq$ constraints on independent scalar Gaussian coordinates.
The orthogonal coordinates of each internal row retain their Gaussian
law and remain independent of the constrained coordinates. Thus
\[
 W^s=W_*^s+R^s+E^s,
 \qquad R^s=R^sP,\quad E^sP=0,
 \qquad \|R^s\|_{\mathrm{op}}\le\rho t\sqrt C.
\]
Conditionally on the cylinder, each free matrix $E^s$ has law
$(t/\sqrt N)G^s(I-P)$ with $G^s$ a standard Gaussian matrix.

\begin{lemma}[Free Gaussian directions]\label{lem:operator}
Conditionally on $\mathcal C_\rho$, all internal free matrices satisfy
$\|E^s\|_{\mathrm{op}}\le b_Lt$ with probability at least $1/2$.
\end{lemma}
\begin{proof}
A $1/4$-net of the unit sphere in $\R^N$ has at most $9^N$ points,
by the disjoint-ball volume argument. Approximating both vectors in a
bilinear form gives
$\|G\|_{\mathrm{op}}\le2\max_{u,v\text{ in the net}}|u^TGv|$.
Each such bilinear form for standard Gaussian $G$ is standard normal.
The scalar Gaussian tail bound gives
\[
 \PP\{\|G\|_{\mathrm{op}}>b_L\sqrt N\}
 \le2\,9^{2N}e^{-b_L^2N/8}
 =2(4L)^{-N}\le\frac1{2L}.
\]
Use $\|G(I-P)\|_{\mathrm{op}}\le\|G\|_{\mathrm{op}}$ and a union
bound over the $L-2$ internal matrices. For $L=2$ the event is automatic.
\end{proof}

\begin{lemma}[Uniform preservation of the target margin]\label{lem:robustness}
Every draw in $\mathcal C_\rho$ satisfying the event of
Lemma~\ref{lem:operator} implements $f$ with strict margin greater than
$3a/2$ on the entire cube.
\end{lemma}
\begin{proof}
The choice of $D_0$ implies $0<\rho_0<1$, and therefore
$\rho t=\rho_0t^{-(L-2)}\le1$.
For each input define
$e_s=N^{-1/2}\|h_w^s-h_*^s\|_2$.
The first-layer constraints and $\tanh$'s Lipschitz constant give
$e_1\le\rho t(d+1)$.
Since $E^s h_*^{s-1}=0$, unit boundedness of the hidden activations and
\eqref{eq:clone-norms} give
\[
 e_s\le(\sqrt2 C^2+\rho t\sqrt C+b_Lt)e_{s-1}+\rho t\sqrt C
      \le tA e_{s-1}+\rho t\sqrt C.
\]
Summing the geometric series and using $\sqrt C/(tA-1)\le1$ yields
\[
 e_s\le\rho t^s(d+2)A^{s-1}.
\]
In particular,
$e_{L-1}\le\rho_0(d+2)A^{L-2}$.
At the output the perturbation is bounded by
\[
 |z_w-z_*|\le4a\sqrt2\,C^2e_{L-1}+\rho t
 \le\rho_0[4a\sqrt2\,C^2(d+2)A^{L-2}+1]
 =a/2.
\]
These bounds hold simultaneously for every input, without a union bound
on cube points. The teacher has strict margin greater than $2a$.
\end{proof}

\section{The probability cost and posterior bound}
\begin{lemma}[Cylinder probability]\label{lem:cylinder-cost}
The target-mass bound~\eqref{eq:mass} holds with $B_N$ from
\eqref{eq:cost}.
\end{lemma}
\begin{proof}
The sum of squared centers of the $Nq$ constrained normalized coordinates
is at most $NB_0$. A first-layer center has magnitude at most $\sqrt d C$.
An internal center has magnitude at most $\sqrt{2C}\,C$.
A readout center has magnitude at most $8aC^2$.
Counting the coordinates of each kind gives the expression for $B_0$ in
\eqref{eq:constants}.

For a centered Gaussian $Z$ of standard deviation $t$, integrating the
minimum density on an interval and using
$(|c|/t+\rho)^2/2\le c^2/t^2+\rho^2$ gives
\[
 \PP\{|Z-c|\le\rho t\}
 \ge\frac{2\rho}{\sqrt{2\pi}}e^{-c^2/t^2-\rho^2}.
\]
Multiply over the independent constrained coordinates and use $\rho\le1$.
Lemmas~\ref{lem:operator} and~\ref{lem:robustness} imply
\begin{align*}
 -\log\mu_{N,\gamma}(A_{N,f}^a)
 &\le \frac{NB_0}{t^2}
       +Nq\log\frac{e\sqrt{2\pi}}{2\rho}+\log2\\
 &=B_0N^\gamma
   +Nq\left[\log\frac{e\sqrt{2\pi}}{2\rho_0}
       +\frac{(L-1)(1-\gamma)}2\log N\right]+\log2.
\end{align*}
This is $B_N$. Because $N^\gamma\le N$, while $B_0,q$, and
$\rho_0^{-1}$ are polynomially bounded in $d,C$ at fixed $L,a$,
\eqref{eq:order} follows.
\end{proof}

\begin{lemma}[A target-mass posterior guarantee]\label{lem:bayes}
Let $\pi$ be a prior on a finite class of Boolean functions and suppose
$\pi(f)\ge e^{-B}>0$. Under exact consistency conditioning on $n$
independent noiseless examples, the sample bound
\[
 n\ge(6/\varepsilon)[B+\log(6/(\varepsilon\delta))]
\]
implies $R\le\varepsilon/3$ and $F\le\varepsilon$ with probability
at least $1-\delta$.
\end{lemma}
\begin{proof}
Let $Z(S_n)$ be the posterior normalizer. The target is always consistent,
so $Z(S_n)\ge\pi(f)$. A fixed function of true error $r$ survives all
labels with probability $(1-r)^n$. Thus, if $U(S_n)$ denotes posterior
mass on functions of error greater than $\varepsilon/6$,
\[
 \E_{S_n}U(S_n)
 \le\frac1{\pi(f)}
   \sum_{g:\operatorname{err}_D(g)>\varepsilon/6}
       \pi(g)(1-\operatorname{err}_D(g))^n
 \le e^{B-n\varepsilon/6}.
\]
Markov's inequality gives $U\le\varepsilon/6$ with probability
at least $1-\delta$ at the displayed sample size.
The remaining posterior mass has error at most $\varepsilon/6$, so
$R\le\varepsilon/3$. On an input with correct-label probability at most
$2/3$, the posterior error probability is at least $1/3$.
Consequently $F\le3R\le\varepsilon$, including equality in the failure event.
\end{proof}

\begin{proof}[Proof of Theorem~\ref{thm:learning}]
Lemma~\ref{lem:cylinder-cost} proves the raw mass bound and its growth
estimate. The global-margin normalizer is at most one, so
$\pi_N(f)\ge\mu_{N,\gamma}(A_{N,f}^a)\ge e^{-B_N}$.
In particular all posterior normalizers are positive.
Apply Lemma~\ref{lem:bayes}.
\end{proof}

\section{Parity and simultaneous growth}
\begin{lemma}[A small parity network]\label{lem:parity}
At every fixed depth $L\ge2$, full parity on $k$ bits has
$C_L(f)=O_L(k\log(k+1))$.
\end{lemma}
\begin{proof}
Put
\[
 t(x)=\frac{k+\sum_{i=1}^k x_i}{2},\qquad
 y_j=(-1)^{k-j},\qquad A_k=\log(16(k+1)).
\]
At integer $t\in\{0,\ldots,k\}$, define
$s_j(t)=\operatorname{sign}(t-j+1/2)$ for $1\le j\le k$.
The identity
\[
 y_t=c+\sum_{j=1}^k y_js_j(t),\qquad c=(y_0+y_k)/2\in\{0,1\},
\]
is checked by its value at $t=0$ and its jumps.
Replace $s_j$ by $h_j=\tanh(A_k(t-j+1/2))$.
Each replacement has error at most $2e^{-A_k}$, since its argument
has absolute value at least $A_k/2$. A constant neuron $h_0=\tanh1$
provides the term $c$ with readout weight $c/\tanh1$.
The resulting output differs from parity by less than $1/8$.
This network has width $k+1$, first-layer weights bounded by
$O(k\log(k+1))$, and readout weights bounded by two.

For $L\ge3$, let the next hidden layer apply $\tanh$ to this sum.
Its correct-signed output is greater than $\tanh(7/8)>1/2$.
Every additional hidden layer preserves the margin using a single
weight two, and the final readout has weight one. Pad unused neurons
with zero weights. Finally pad the common width to an integer
$C=O_L(k\log(k+1))$ that also bounds every weight. Since $\tanh0=0$,
this preserves the implementation.
\end{proof}

\begin{proof}[Proof of Corollary~\ref{cor:joint}]
The polynomial-width statement follows directly from
\eqref{eq:order} and~\eqref{eq:samples}. For parity use
Lemma~\ref{lem:parity}. At fixed accuracy and confidence parameters, the
sufficient sample size is bounded by a polynomial in $k$ times
$N(k)\log(eN(k))$. Its logarithm is $o(k)$ whenever
$\log N(k)=o(k)$.
\end{proof}

The cylinder uses $N[d+2+(L-2)C]$ constrained directions even though a
dense deep network has quadratically many parameters in width. The other
internal directions are retained at their prior scale; their effect is
controlled by shrinking the constrained radius. For $\gamma<1$, this
shrinking costs an additional $\log N$ per constrained direction.
That finite-width cost is why these learning guarantees coexist with
possible statistical obstructions after taking width to infinity first.
\end{document}
